%% kapa: exam=1 round=25 year=2014 period=1 form=A subject=civil_law subject_ko=민법 seq=38 number=38
%% source: Q-Net 감정평가사 기출문제 (cst00309) artlSeq=5101207 file=extracted/1008/5101207/2014년 제25회 감정평가사 1차시험 기출문제/1교시(A형).hwp
%% conversion: kordoc 4.24.51 -> tools/gen1.py (auto); review=proofread-sonnet; flags=none
\begin{question}{38}
저당권의 효력에 관한 설명으로 옳지 않은 것은? (다툼이 있으면 판례에 의함)
\begin{choices}
  \item 저당부동산에 관하여 이해관계인이 없는 경우, 원본의 이행기를 경과한 후의 1년분 이상의 지연손해에 대해서도 저당권의 효력이 미칠 수 있다.
  \item 건물에 저당권이 설정된 경우, 원칙적으로 그 건물에 부속된 창고에도 저당권의 효력이 미친다.
  \item 저당권의 효력이 미치는 종물은 저당권 설정 전부터 존재하였던 것이어야 한다.
  \item 건물에 대한 저당권의 효력은 원칙적으로 그 대지이용권인 지상권에도 미친다.
  \item 저당권과 전세권이 경합하는 경우에는 설정등기의 선후에 의하여 우선순위를 정한다.
\end{choices}
\end{question}
\answer{3}
